% ------------------------------------------------------------
% Chapter 5: Discussion
%
% Results says what happened. This chapter says what it means and what it does
% not mean. Every number is reproduced from Chapter 4 and traces to a frozen
% table under tables/; nothing is recomputed in prose.
%
% The eight divisions are \paragraph and carry \label{sec:...}, so none of them
% reaches the table of contents. Promote them to \section if an examiner should
% be able to navigate the chapter.
% ------------------------------------------------------------

\chapter{Discussion}
\label{ch:discussion}

\paragraph{What the Experiment Establishes}
\label{sec:discussion-establishes}

Because the scenario, the phrasing, the sampling parameters and the replicate
count are constant across conditions, a difference between two conditions is
attributable to the age signal and to nothing else the experiment manipulated.
Because the age is supplied rather than inferred from a real interaction
history, no difference is attributable to a model having worked out who it is
talking to.

A model that behaves differently towards a user who has written that they are
twelve has demonstrated age conditioning, not age recognition, and the two are
separate capabilities with separate failure modes. A deployed service that
verifies age at the account layer and passes a verified age to the model
exercises the first; a service that infers age from conversational content
exercises the second, and this study says nothing about how well it would do
so. The implicit cue conditions probe that boundary, and what they establish is
behavioural: a cue moves the reply less than a statement does, and the design
cannot say whether the cue was less often noticed or noticed and weighted less.
Section~\ref{sec:discussion-signals} returns to this.

\paragraph{Threshold and Gradient}
\label{sec:discussion-shape}

The central claim of this thesis is not that models treat minors differently.
That was expected and would not have needed a benchmark of this size. The claim
is that the differential treatment has a different shape on different dimensions
of the same reply.

On safety, the conditioning is concentrated. Refusal on age-restricted requests
falls by 27.8 percentage points in the single year from seventeen to eighteen,
which is 48.3\% of the entire movement across the fourteen years the latter
spans. On linguistic accessibility, it is distributed: the same step accounts for
between 4.0 and 11.0\% of each model's full range, and on two of the six its
interval includes zero. One outcome behaves as though a rule is applied at a
boundary, the other as though a continuous adjustment is made to an estimate of
the reader.

Two readings are available, and the data distinguish them only partly. Either the
two behaviours have different origins, a safety boundary reflecting an explicit
policy encoded during alignment and a linguistic adjustment reflecting a learned
regularity about how people write for children, or both are gradients and the
safety gradient is much steeper near eighteen. A pure gradient would not
ordinarily place half its movement in one of fourteen years, and no adjacent
step within the minor ages survives correction whereas the boundary step does on
five models. But a behavioural experiment cannot see a policy, only its shadow,
Moreover, a sufficiently sharp gradient is observationally close to a threshold.

Mistral Small 4 is the instructive exception. It conditions on age about as
strongly as the panel when minor and adult ages are compared as blocks, at 22.9
points, and barely at all across the year the statutes actually name, at 2.7
points with an interval crossing zero. That is age conditioning without a
boundary, and it shows that a model can adjust to age and still fail what an
age-assurance regime is trying to achieve, which is different treatment either
side of a legal line.

\paragraph{Refusal, Delivery, and Harm}
\label{sec:discussion-delivery}

The four-cell outcome is the methodological contribution most likely to
transfer, and its value is clearest in the one model where it changes the
answer. Read on refusal alone, Mistral Small 4 is the most permissive model on
the panel at 17.5\%; read on delivery alone, it is the least useful at 65.8. Both
readings are wrong, because 16.9\% of its replies agree to help and then supply
nothing. A binary Refusal Rate, which is what most published safety evaluations
report, scores those replies as compliance and places that model at the
permissive end of a ranking it does not belong at.

The decomposition should not be oversold in the other direction. The cell it was
built to expose, the reply that declines and then answers anyway, turned out to
be rare, at 0.2\% on the macro-average and never above half a point on any
model. The mismatch that turned out to matter is the opposite one, and it was
not the one anticipated.

A more serious limit runs through the whole safety analysis. Delivery Response
records that the material was supplied, not whether it was harmful, so a reply
answering a restricted question for a fourteen-year-old with a careful, hedged,
partial answer and one answering it with operational detail are scored
identically. The directional diagnostic is named for what it counts: delivery
where a refusal was expected, and stands as an upper bound: the macro-average of
30.6\% on harmful requests is not a claim that three in ten replies harmed
anybody. Closing that gap needs a severity judgement the rubric does not make,
and Chapter~\ref{ch:future} sets out the audit that would supply one.

A second limit is of a different kind, and Section~\ref{sec:semantics} makes it
visible. The four-cell outcome and the thirteen rubric labels are closed
vocabularies, so they register only the changes they were built to name, and a
reply can be rewritten in register, framing, and what it offers and still occupy
the same cell. Introducing an explicit minor age where Neutral gave none
separates the text by 0.189 $[0.164, 0.216]$ on age-restricted requests, the
single step from seventeen to eighteen separates it by 0.079 $[0.058, 0.104]$
where the same step moves the three control strata by between 0.018 and 0.029,
and the association between that displacement and whether the classified outcome
changed is never above $\rho = 0.202$.

The measure has no polarity: a reply rewritten into gentler framing with a
signpost to a trusted adult and one rewritten into evasion are equally
displaced. It therefore does not extend the safety claim and is not a second
measurement of harm. It bears on scope. Age conditioning here is broader than
any categorical measure in this thesis records, so the rates of
Section~\ref{sec:safety} describe one visible face of the behaviour, and an
evaluation built only on category counts will understate how much a disclosed
age changes.

\paragraph{Explicit and Implicit Age Signals}
\label{sec:discussion-signals}

The ordering of the signal levels is one of the cleaner results in the thesis
and one of the easiest to over-interpret. Refusal on age-restricted requests
sits at 13.5\% under Neutral, 12.9 with an adult-associated cue, 16.7 with an
explicit adult age, 29.5 with a minor-associated cue and 62.0 with an explicit
minor age. The two adult-side conditions are indistinguishable from Neutral,
which is what one would expect if the default assumption is already an adult
user, and the minor cue moves roughly a third of the way to the explicit minor
age. The pattern repeats on reading level.

What cannot be said is that the cue was received and then discounted. That
describes a two-stage process, recognition followed by weighting, and the
experiment observes only the composite: a model recognising a mention of a
school timetable as a minority signal on one occasion in three and conditioning
fully when it does, and a model recognising it every time and conditioning at a
third of the strength, produce the same aggregate. Separating them requires
eliciting the model's belief about the user rather than only its behaviour.

The practical implication survives the ambiguity. Real users rarely announce
their age; they reveal it through what they talk about and how they write, which
is the implicit condition. If conditioning on a cue is a third of conditioning
on a statement, the protection a model offers unprompted is materially weaker
than the protection it offers when told, and an evaluation that only ever states
the age reports the stronger number.

\paragraph{Heterogeneity Across the Panel}
\label{sec:discussion-panel}

Six models from five providers do not behave as one model with noise around it,
and the disagreement is structured. Gemma 4 31B is strongest on the adult side
of the age-restricted stratum at 92.7\% alignment and among the weakest on the
minor side at 55.1, applying the boundary more readily in the permissive
direction than the protective one. GPT-5.6 Luna inverts on the two directional
diagnostics, worst on the panel for delivery on harmful requests at 36.5\% and
best for delivery to an explicit minor age at 19.1. Gemini 3.5 Flash Lite is the
only model whose requests are withheld at the provider layer, and that happens
exclusively at minor ages.

That last case is a different mechanism reaching a similar end. Seventy-nine
age-restricted requests were withheld before the model saw them, all at explicit
minor ages and concentrated at the younger end, and counted end to end that
model's delivery to an explicit minor age falls from 31.2 to 26.7 per cent. A
benchmark scoring only returned replies credits the model for restraint it did
not exercise, and one scoring only end-to-end cannot distinguish a model that
declined from a request that never reached one. Reporting both is the only
defensible option, and the practical lesson for benchmark design is that the
provider stack is part of the model under test.

A macro-average over six models describes this panel and does not estimate a
population of models, because the panel is fixed rather than sampled. Where the
six disagree in direction, as at the statutory boundary, the per-model column is
the result.

\paragraph{Developmental Accessibility as an Operational Construct}
\label{sec:discussion-accessibility}

The linguistic analysis rests on a mapping from an age to a target grade, and
its honesty depends on that mapping being operational rather than normative. The
target is a calibration reference derived from a schooling convention, not a
claim about what a particular child can read, which varies enormously within any
age, nor about what a model ought to produce. Age-Alignment Error is an error
against a stated reference and not a developmental deficiency, so the finding
that replies sit 2.31 to 2.94 grades from it is a statement about calibration
rather than about harm.

Within those limits, the result is still meaningful. Reading level moves in the
right direction on all six models and across every floor threshold tested, and
it reaches the reference on none of them. A model that adapts directionally but
under-adapts in magnitude is a recognisable engineering condition and a
tractable one.

The decomposition corrects a natural assumption. The reduction in grade level is
carried almost entirely by the syllables-per-word term, at $-1.80$ grades, while
the sentence-length term contributes $+0.03$ with an interval containing zero
and a sign that is not stable across the panel. Models writing for an explicit
minor age choose shorter words; they do not write shorter sentences. The
intuitive account runs the other way, and the two channels differ in what they
do: a lexical simplification changes which concepts are nameable, a syntactic
one how much has to be held in mind at once. Neither has been shown to improve
comprehension, since a readability formula predicts difficulty from surface
features and no comprehension study was conducted here.

\paragraph{Implications for Age Assurance}
\label{sec:discussion-assurance}

The regulatory direction of travel treats knowing the user's age as the
substantive step, with appropriate treatment following from it. The results here
suggest that framing is incomplete in a specific and testable way. Supplying the
age is not the hard part of the problem for these models, which clearly use it.
The hard part is that what follows from the age is uneven across the dimensions
on which a young user is actually exposed.

Three consequences follow, and each is narrow. First, an age-assurance
obligation discharged at the account layer will produce a large behavioural
change at the statutory boundary and a much smaller one across childhood itself,
so a regime concerned with a fifteen-year-old rather than with the seventeen to
eighteen transition is relying on the part of the response surface that moves
least. Second, an evaluation regime that certifies a model on explicit age will
overstate the protection available to users who never state one, by roughly a
factor of three on the evidence here. Third, a compliance test built on Refusal
Rates alone can be satisfied by a model that agrees and then supplies nothing,
which is neither safe nor useful, and the four-cell outcome or something like it
is needed to tell the two apart.

None of this bears on whether age assurance is desirable, which is a policy
question this thesis does not address, and none of it establishes that any model
tested breached any obligation. What it supports is narrower: verifying age and
adapting to age are separate problems, and progress on the first does not entail
progress on the second.

\paragraph{Strengths and Weaknesses of the Design}
\label{sec:discussion-design}

The design's principal strength is that it isolates the variable of interest.
Two hundred scenarios crossed with thirteen conditions and three replicates give
46,800 requests in which everything but the age signal is held fixed, and the
pairing is within scenario, so a contrast compares a scenario against itself.
The inferential machinery is declared before the analysis: families are
registered, tests are permutation-based on the pairing the design creates,
intervals resample scenarios rather than replies, planned controls are read
against a pre-registered margin rather than a failure to reject, and the
classifier is validated with an admission threshold fixed in advance. That
covers Sections~\ref{sec:safety} to~\ref{sec:dialogue} and not
Section~\ref{sec:semantics}, which is the last weakness below.

Six weaknesses are material, and none is fully remediable within this study.

The first is that the classifier is itself a model. Agreement with the annotator
is strong on the two primary labels of the single-turn corpus and weak on
several characteristics, and the weak ones are reported descriptively for that
reason. On the dialogue arm, the second primary label does not hold:
Section~\ref{subsec:judge} reports Delivery Response below the admission
threshold at both pressed turns, so the measures of Section~\ref{sec:dialogue}
that read the delivery axis are the weaker ones. Agreement is also not accuracy,
since a systematic error shared by classifier and annotator would not appear in
a $\kappa$. The shared-provider concern is answered as far as it can be, by the
classifier disagreeing most on the model from its own provider.

The second is the length floor, set out in
Section~\ref{subsec:readability-floor}. It removes 9.6\% of returned replies
unevenly, from 1.1\% of DeepSeek-V4 Flash to 19.9\% of Mistral Small 4, so the
readability sample is not the response corpus and differs from it in a way
correlated with both the model and the stratum. No sensitivity analysis can
reconstruct the reading level of a reply too short to have one. The model losing
most replies to the floor is the one that agrees and then supplies nothing, so
the two findings are partly the same phenomenon seen twice.

The third is that the main corpus is single-turn, English and synthetic. That is
why Experiment~2 exists, and its results are now in hand:
Section~\ref{sec:dialogue} reports the age gap on age-restricted requests
eroding from 41.3 points at the opening to 18.5 by the third turn, while Action
Defect rises by 23.1 points over the same two turns. That addresses the
limitation without removing it, since those results are descriptive rather than
tested, cover fifty of the two hundred scenarios and two pressed turns rather
than a sustained conversation, and the delivery axis several of them read does
not clear the admission threshold at depth. The scenarios are also authored
rather than observed, so they represent a considered view of what a young user
might ask rather than a sample of what young users do ask.

The fourth is that the design supplies the age rather than eliciting it, so
recognition is never measured on the same replies. An absence of behavioural
change under a cue condition cannot therefore be attributed to a failure to act
on the signal rather than to a failure to register it. The cue conditions are
read as the weaker arm for this reason among others, and the explicit ages carry
any claim that depends on the signal having been received.

The fifth is temporal. Six commercial models were tested over a fixed window
against provider stacks that change without notice, so a provider shipping an
age-assurance feature next month invalidates the specific numbers while leaving
the method intact. That is a condition of empirical work on deployed models
rather than a defect of this one, but it does mean the contribution is better
understood as a reusable instrument than as a permanent measurement.

The sixth is that Section~\ref{sec:semantics} was specified after the results it
sits beside had been read, and it is labelled exploratory throughout. Three
things limit the damage without repairing it: the measure has no polarity and so
cannot be tuned toward a favourable direction, its three control strata were
fixed by the benchmark rather than chosen for this analysis, and a second
encoder was pinned with the statements that would count as reproduction written
down before any full-corpus result was seen. Cosine similarity also has no
interpretable scale, so comparisons are between strata and never against an
absolute value. None of that makes the section confirmatory. A pre-registered
replication on a fresh scenario draw would.